\section{Wet and dry soil conditions}
\label{sec:moisture}

We used weekly SMAP L4 surface soil moisture to describe departures from
usual seasonal conditions during April--July 2019 and June--August 2022,
the two event windows selected for the flood and drought case studies.
The analysis covered cells eligible for the crop-history analysis and
grouped them by state and event-year CDL label. The crop strata were corn
(1), soybean (5), winter wheat/soybean double crop (26), and oats (28).
SMAP retains its native 9~km spatial support when assigned to the finer
crop grid, so neighboring crop cells can share the same moisture
information.

For each cell and calendar week (ISO week), we estimated a reference mean
and sample standard deviation from 2015--2021, giving at most seven annual
observations. We expressed each event observation as a departure from that
mean in standard-deviation units, with stored $z$-scores clipped to
$[-5,5]$. The reference includes 2019 but precedes 2022. Including 2019
can dampen its anomaly. We treat these as retrospective case studies;
holding out both events would require excluding each event year from its
reference.

We also fitted an empirical-Bayes Normal--Inverse-Gamma model to account
for uncertainty in the reference mean and variance. Its conjugate update
gives a Student-$t$ posterior predictive distribution \citep{murphy2007conjugate}.
The prior mean was the regional mean for the corresponding week; the
variance prior used the regional median of the cell-level baseline
variances. We evaluated the predictive cumulative distribution at each
observed moisture value. This moisture
percentile approaches zero for unusually dry observations and one for
unusually wet observations. It is a position within the fitted moisture
distribution, not a probability that drought is present.

Figure~\ref{fig:moisture} compares mean moisture percentiles for corn and
soybean in each state, with a reference line at 0.5 on the zero-to-one
scale. We averaged the individual pixel-week scores within each
state--crop group and also counted fractions beyond selected thresholds.
A cell observed in several weeks contributes several times. These
fractions therefore describe the combined spatial and temporal extent of
unusual conditions within the sampled population.

The 2019 mean moisture percentiles exceeded 0.5 for corn and soybean in
every study state, reaching 0.85 for South Dakota corn. In Iowa, 17.75\% of
2,901,150 corn pixel-weeks exceeded $z=1.5$, indicating repeated or
widespread departures above the weekly reference. These elevated values
describe relative wetness; identifying flooded fields would require
separate inundation observations.

The 2022 window showed a different pattern. Nebraska soybean had a mean
moisture percentile of 0.24, and 40.54\% of its 947,648 pixel-weeks fell
below the tenth predictive percentile. North Dakota corn averaged 0.52,
showing the geographic variation in the dry signal. Kentucky's winter
wheat/soybean double-crop stratum had 41.96\% of 49,543 pixel-weeks below
the tenth percentile.
Differences among these crop strata reflect their spatial distributions
as well as the moisture conditions encountered, which limits comparisons
of crop-specific responses to water stress.

These summaries identify where and when surface moisture was unusual
within the sampled crop-history population. Read alongside seasonal NDVI,
they provide context for crop-condition monitoring.
